%% ma: L12-B06-0009
%% lop: 12
%% chuong: 2
%% bai: 6
%% dang: 12.6.3
%% loai: TN4
%% muc: TH
%% dap_an: "B"
%% nguon: "(NBV)-ĐỀ SỐ 12-MỖI NGÀY 1 ĐỀ THI 2025.pdf (D12, MỖI NGÀY 1 ĐỀ - Đề số 12), Phần I Câu 10"
%% kien_thuc: "Bài 6 (tổng, hiệu hai vectơ trong không gian, quy tắc ba điểm)"
%% trang_thai: cho-duyet
%% ly_do: "Đề gốc có hai khẳng định sai (B và D) nên không còn đúng một đáp án. Đã sửa phương án D thành AB + BC = AD - CD (vectơ; dấu trừ giữa AB và BC ở gốc nhiều khả năng là lỗi đánh máy); cần giáo viên xác nhận."
%% ghi_chu: "Đáp án gốc B. Kiểm bằng vectơ ngẫu nhiên: A đúng, B sai, C đúng, D (nguyên văn, dấu trừ) sai. Đã sửa D thành dấu cộng để chỉ còn B sai."
\begin{ex}
Cho tứ diện $ABCD$. Khẳng định nào sau đây sai?
\choice{$\overrightarrow{AB}-\overrightarrow{AC}=\overrightarrow{CB}$}{\True $\overrightarrow{AB}+\overrightarrow{AC}=\overrightarrow{AD}$}{$\overrightarrow{AB}+\overrightarrow{CD}=\overrightarrow{AD}+\overrightarrow{CB}$}{$\overrightarrow{AB}+\overrightarrow{BC}=\overrightarrow{AD}-\overrightarrow{CD}$}
\loigiai{A đúng: $\overrightarrow{AB}-\overrightarrow{AC}=\overrightarrow{CB}$ (quy tắc hiệu).

C đúng: $\overrightarrow{AB}-\overrightarrow{AD}=\overrightarrow{DB}=\overrightarrow{CB}-\overrightarrow{CD}$, suy ra $\overrightarrow{AB}+\overrightarrow{CD}=\overrightarrow{AD}+\overrightarrow{CB}$.

D đúng: $\overrightarrow{AB}+\overrightarrow{BC}=\overrightarrow{AC}$ và $\overrightarrow{AD}-\overrightarrow{CD}=\overrightarrow{AD}+\overrightarrow{DC}=\overrightarrow{AC}$.

B sai: $\overrightarrow{AB}+\overrightarrow{AC}=\overrightarrow{AD}$ chỉ đúng khi $D$ là đỉnh thứ tư của hình bình hành $ABDC$, mà $ABCD$ là tứ diện nên $A,B,C,D$ không đồng phẳng. Chọn B.}
\end{ex}
